\documentclass[10pt,a4paper]{amsart}
\usepackage[margin=19mm]{geometry}
\usepackage{amsmath,amssymb,mathtools}
\usepackage[scr=rsfs,cal=euler]{mathalfa}
\usepackage{xfrac}
\usepackage[unicode,hidelinks,bookmarks=false]{hyperref}
\newcommand{\ie}{i.e.\ }
\newcommand{\cf}{cf.\ }
\newcommand{\resp}{respectively\ }
\newcommand{\Subsection}{\subsection}
\providecommand{\nobreakdash}{\nobreak\hbox{-}}
\setlength{\emergencystretch}{2em}
\multlinegap0pt
\usepackage{bm}
\usepackage{bbm}
\usepackage{dsfont}
\usepackage{xspace}
\usepackage{cancel}
\usepackage{mathtools}
\usepackage{amsthm}
\makeatletter
\@ifundefined{thm}{\newtheorem{thm}{Thm}}{}
\@ifundefined{prop}{\newtheorem{prop}[thm]{Proposition}}{}
\makeatother
\title[Linear strands of powers of certain binomial edge ideals]{Linear strands of powers of certain binomial edge ideals}
\author{Abbas Dohadwala, Bryan Flores-Silva, Alicia Orozco-Moya, Zoe Siegelnickel}
\date{Source: arXiv:2601.10842v1 (CC BY 4.0)}
\begin{document}
\maketitle

\noindent\emph{Source and licence.} Linear strands of powers of certain binomial edge ideals. Version arXiv:2601.10842v1 (CC BY 4.0), distributed under the Creative Commons Attribution 4.0 International licence, as recorded in the arXiv licence metadata for this exact version and shown on the abstract page https://arxiv.org/abs/2601.10842v1. Full rights notice: licenses/arxiv-2601.10842-CC-BY-4.0.txt.\par\smallskip

\noindent\textit{Editorial scope.} Counted unit: the Prop whose source label is \texttt{lem: crucial iso} in the article (source0.tex lines 369--372, source label \texttt{lem: crucial iso}), delivered together with the article's own complete original proof of it (source0.tex lines 374--437, one \texttt{proof} environment of 64 lines). No mathematics is written, repaired or re-derived: statement and proof are the article's own text line for line. Required context retained uncounted: eq: L (lines 357--360). Every definition, display and label of those lines is retained unchanged. Omitted from this package and not redistributed: 1--356, 361--368, 373--373, 438--563. Nothing inside a retained excerpt is omitted or altered: every retained body is delivered exactly as the article's lines are written, so no omission receipt of any kind accompanies this package. Citations: the retained excerpts carry no citation command at all, so no bibliography entry of the article is transcribed and no reference is redistributed. Reference adaptation: none was needed and none was made; every reference inside the delivered unit resolves inside it, so no unresolved reference or citation is produced. Printed slips kept literal: no printed word, symbol, hypothesis or number of the counted statement or of its proof was altered. Layout and class: the article's own class is \texttt{amsart} with packages including geometry, graphicx, amsfonts, amssymb, amsthm, amsmath, tikz, caption, tcolorbox, hyperref, cleveref, enumitem, mathtools; the wrapper is the lane's pinned \texttt{amsart} editorial wrapper at 10pt on A4, which restates the theorem-like environments the retained text needs and adds the article's own macro definitions that the retained text needs (none), with extra wrapper packages (bm, bbm, dsfont, xspace, cancel, mathtools). The delivered numbering is the unit's own. Assets: no retained span contains a figure environment, a graphics-inclusion command, a PDF-page inclusion, a table or a TikZ picture; no image file is copied into this package, the package declares no asset dependency and no image file is redistributed with it. Licence: version arXiv:2601.10842v1 of this article is distributed under the Creative Commons Attribution 4.0 International licence, as recorded in the arXiv licence metadata for this exact version and shown on the abstract page https://arxiv.org/abs/2601.10842v1. A full rights notice is delivered at licenses/arxiv-2601.10842-CC-BY-4.0.txt. Characters: every retained excerpt is pure ASCII, so the delivered engine, XeLaTeX with its default Latin Modern text font, reports no missing character.
\begin{eqnarray}
L_{(m,-)}    & =&\{F\in R[T_1, \ldots, T_s]:\varphi(F)=0, \deg{F}=(m,-)\} \nonumber\\
    & = &\left\{F=\sum_{\alpha\in\mathbb{Z}_{\geq 0}^s, |\alpha|=m} r_\alpha T_{i}^{\alpha_{i}}\in R[T_1,\ldots, T_s] : \sum_{\alpha\in\mathbb{Z}_{\geq 0}^s, |\alpha|=m}r_\alpha f_{i}^{\alpha_{i}}=0\right\}. \label{eq: L}
\end{eqnarray}

\begin{prop} \label{lem: crucial iso}
    Let $R$ be a graded $\mathbb{K}$-algebra  with $R_0=\mathbb{K}$ and let $I$ be a homogeneous ideal of $R$ with minimal homogeneous generating set $\{f_1, \dots, f_s\}$ of elements of degree $d$.
    Then there is an isomorphism between \(\Omega_1(I^m)\) and \(N =L_{(m,-)}/L_{(m,\geq md)}\) as graded $R$-modules.  
\end{prop}

\begin{proof}
    Let \(\{h_1, ..., h_t\}\subset\{ f^\alpha:|\alpha|=m \}\) be a minimal generating set of \(I^m\) as a \(\mathbb{K}\)-vector space. Considering a basis $e_1, \ldots, e_t$ for the free module $R^t$ and elements $r_1,\ldots, r_t\in R$ one has an $R$-module homomorphism
    \[\partial:R^t\rightarrow R 
    , \quad \partial \left(\sum_{i=1}^t r_ie_i\right)= \sum_{i=1}^t  r_ih_i.
    \] 
    Recall that \(\Omega_1(I^m)=\ker(\partial)\).

    

    For each $1\leq i\leq t$, fix  $\alpha^{(i)}=(\alpha^{(i)}_1, \ldots, \alpha^{(i)}_s)\in \mathbb{Z}_{\geq 0}^s$ such that the minimal generator $h_i$ decomposes as \(h_i=f_1^{\alpha^{(i)}_1}\cdots f_s^{\alpha^{(i)}_s}\).
    Consider the following map
    \begin{equation}\label{eq: defg0}
    g_0:\Omega_1(I^m)\rightarrow R[T], \quad g_0\left(\sum_{i=1}^t r_ie_i\right)=\sum_{i=1}^t r_iT_1^{\alpha_1^{(i)}}\cdots T_s^{\alpha_s^{(i)}}.
    \end{equation}
    
    Applying \(\varphi\) to an element in the image of $g_0$ yields 
    \begin{equation*}
        \varphi\left(\sum_{i=1}^t r_iT_1^{\alpha_1^{(i)}}\cdots T_s^{\alpha_s^{(i)}}\right)
        = T^m\sum_{i=1}^t r_ih_i
        = T^m \partial\left(\sum_{i=1}^t r_ie_i\right)= 0.
    \end{equation*}
    Thus, the image of $g_0$ is contained in $L_{(m, -)}$ by \eqref{eq: L}. 
    The map \(g_0\) induces an $R$-module homomorphism \(\Omega_1(I^m)\xrightarrow{g_1} N\) which maps elements of degree $v$ in the source to elements of degree $(m,v)$ in the target, thus preserving $R$-degrees.

 Next we define an inverse mapping. 
  Recall from \eqref{eq: L} that
    \begin{align}\label{eq:observe}
\sum_{\,|\alpha^{(i)}\,|=m} r_iT^{\alpha^{(i)}} = &\sum_{\,|\alpha^{(i)}\,|=m} r_iT_1^{\alpha_1^{(i)}}\cdots T_s^{\alpha_s^{(i)}}\in L_{(m,-)} 
        \iff& \sum_{\,|\alpha^{(i)}\,|=m} r_i f_1^{\alpha_1^{(i)}}\cdots f_s^{\alpha_s^{(i)}}=0.
    \end{align}
Since \(h_1,...,h_t\) generate  \(I^m\) and \(f_1^{\alpha_1^{(i)}}\cdots f_s^{\alpha_s^{(i)}}\in I^m\) whenever $|\alpha^{(i)}\,|=m$,  we can write   
\begin{equation}\label{eq:aijdef}
f_1^{\alpha_1^{(i)}}\cdots f_s^{\alpha_s^{(i)}}=\sum_{j=1}^t a_{ij}h_{j},
\end{equation}
for some \(a_{ij}\in\mathbb{K}\). 
    Then equation \eqref{eq:observe} becomes
    \begin{equation}\label{eq: becomes}
        \sum_{\,|\alpha^{(i)}\,|=m} r_iT^{\alpha^{(i)}} \in L_{(m,-)}\iff \sum_{\,|\alpha^{(i)}\,|=m} r_i \left(\sum_{j=1}^ta_{ij}h_{j}\right)=0.
        \end{equation}
    Using notation introduced in \eqref{eq:aijdef}, we define
  
    \begin{equation}\label{eq: defP0}
        p_0:L_{(m,-)}\rightarrow \Omega_1\left(I^m\right),\quad
        p_0\left(\sum_{|\alpha^{(i)}|=m} r_i T_1^{\alpha_1^{(i)}}\cdots T_s^{\alpha_s^{(i)}}\right)=\sum_{\,|\alpha^{(i)}\,|=m} r_i \left(\sum_{j=1}^t a_{ij} e_{i}\right).
    \end{equation}
    By \eqref{eq: becomes} it follows that
    \begin{align*}
        \partial\left( \sum_{\,|\alpha^{(i)}\,|=m} r_i \sum_{j=1}^t a_{ij} e_{j}\right)
        =&\sum_{\,|\alpha^{(i)}\,|=m} r_i \left(\sum_{j=1}^t a_{ij} h_{j}\right)=0.
    \end{align*}
    Thus the image of $p_0$ is contained in the syzygy module \(\Omega^1(I^m)\).

    We now wish to show that $p_0: L_{(m, -)} \to \Omega_1(I^m)$ induces a well defined map $\overline{p_0}: N \to \Omega_1(I^m)$. 
    It is enough to show that $L_{(m,md)}$ is contained in the kernel of $p_0$. For \(F=\sum_{|\alpha^{(i)}|=m}c_iT_1^{\alpha_1^{(i)}}\cdots T_s^{\alpha_s^{(i)}}\) we see that \(\deg{(F)}=(m,md)\) if and only if $\deg_R(c_i)=0$ if and only if $c_i \in \mathbb{K}$. 
    Then setting  \(b_j=\sum_{i} c_ia_{ij}\in\mathbb{K}\) yields
    \[p_0(F)=p_0\left(\sum_{i}c_iT_1^{\alpha_1^{(i)}}\cdots T_s^{\alpha_s^{(i)}}\right)=\sum_{i}c_i\left(\sum_j a_{ij}e_{j }\right)=\sum_{j=1}^t b_je_j.\]
   
    If $F\in L_{(m,mn)}$ then \(p_0(F)=\sum_{j=1}^t b_je_j\in\Omega_1(I^m)\), thus we have $\partial\left(\sum_{j=1}^t b_je_j \right)=\sum_{j=1}^t b_jh_j=0$.
 
    It follows by linear independence of $h_1, \ldots, h_s$ over $\mathbb{K}$ that \(b_j=0\) for all $1\leq j\leq t$. Therefore, \[p_0(F)=\sum_{j=1}^t b_je_j=0.\]
    
    We have shown that \(L_{(m,md)}\) is contained in the kernel of \(p_0\), thus  there is a well-defined $R$-module homomorphism $\overline{p_0}:N\to \Omega_1(I^m)$. 
    Lastly one can check that $p_0$ and $g_1$ are mutually inverse maps using equations \eqref{eq: defg0} and \eqref{eq: defP0}, and the fact that $g_1$ is induced by $g_0$.
\end{proof}

\end{document}
